% https://web.archive.org/web/20150413003534/http://www.ee.nmt.edu/~teare/ee308l/datasheets/S12SPIV3.pdf
% https://www.bitsavers.org/components/motorola/_appNotes/AN-0991_Using_the_Serial_Peripheral_Interface_to_Communicate_between_Multiple_Microcomputers.pdf
% https://onlinedocs.microchip.com/oxy/GUID-835917AF-E521-4046-AD59-DCB458EB8466-en-US-1/GUID-E4682943-46B9-4A20-A62C-33E8FD3343A3.html
% https://onlinedocs.microchip.com/oxy/GUID-199548F4-607C-436B-80C7-E4F280C1CAD2-en-US-1/GUID-E0991B55-E5C3-419C-AD24-6746DD6BA228.html
% Bitácora: Raíces primitivas y orden multiplicativo
\documentclass[12pt]{article}
\usepackage{algorithm}
\usepackage{algorithmic}
\input{generalconfig.tex}
\title{Communication Protocols}
\author{Contreras Reyes Orlando}
\date{\today}
\begin{document}
\maketitle
\begin{abstract}
This document aims to explain the SPI protocol from the scratch and aspires to replicate the protocol itself in hardware with its four modes.
\end{abstract}

\section{Introduction}
To understand the value of the SPI protocol, it is neccesary to un
\newpage
\tableofcontents
\newpage
The functional block of the SPI is the next

Here it shows the Configuration signals \textbf{CPOL} and \textbf{CPHA}, the main clock (CLK) and the data loaded from ROM (DATAIN). 
Also we could appreciate the \textbf{MOSI} and \textbf{MISO} signals, which are the data lines for the main and secondary devices\footnote{It was decided to keep the old convention of the signals, but changing Master to Main and Slave to Secondary}. The \textbf{SS} signal is used to select the secondary device that will communicate with the master. The \textbf{SCLK} signal is the clock signal generated by the main device to synchronize the data transfer.

Considering the official document of SPI by Motorola
%https://www.bitsavers.org/components/motorola/_appNotes/AN-0991_Using_the_Serial_Peripheral_Interface_to_Communicate_between_Multiple_Microcomputers.pdf
First we need to start with the main enable of the system, the \textbf{NSS}(Secondary Select), we could see that it's a \textbf{active low} signal, when is low the main device is enabled and the communication can start, when is high the main device is disabled and the communication is stopped. So we define an internal signal for the SPI block called \textit{SPI\_EN} we define as the negated of NSS
\begin{verbatim}
SPI_EN <= not NSS;
\end{verbatim}

We could observe on the Data Clock Timing Diagram the behavioral of CPOL and CPHA, we could appreciate that the CPHA signal determines the clock edge on which the data is sampled, while the \textbf{CPOL} signal determines the \textbf{idle state of the clock}. 

The combination of these two signals allows for four different modes of operation, which can be selected based on the requirements of the communication.

Here it appears and important concept \textbf{Leading and Trailing edge}

We define the Leading Edge as the back of the signal and Trailing as the front like if we were talking about an arrow. Therefore if we had CPHA = 0 we select the Leading Edge, then if we set CPHA = 1 we select the Trailing Edge.

The Rising Edge would be on the next cases
$$
\begin{cases}
\text{CPOL = 0, CPHA = 1} \\
\text{CPOL = 1, CPHA = 0}
\end{cases}
$$
Therefore the Rising is determine by a XOR between CPOL and CPHA


The Falling Edge would be on the follow cases
$$
\begin{cases}
\text{CPOL = 0, CPHA = 0} \\
\text{CPOL = 1, CPHA = 1}
\end{cases}
$$

As we previously mention the CPOL signal determines the idle state of the clock, if CPOL = 0 the clock is low when idle, and if CPOL = 1 the clock is high when idle. With this we have the next behavioral
$$
\begin{aligned}
\begin{cases}
    Oscilating &\text{ when } spi_{en} = 1\\
    CPOL       &\text{ when } spi_{en} = 0
\end{cases}
\end{aligned}
$$
We could built this with a procces block\footnote{For Sequential logic is highly reccomended to use procces to assure the correct behavior of the circuit, and to avoid glitches.} 
\begin{verbatim}
process(clk)
begin
if rising_edge(clk) then
        if spi_en = '1' then
                sclk_out <= not sclk_out;
        else
                sclk_out <= CPOL;  -- idle state
        end if;
end if;
end process;
\end{verbatim}


\begin{thebibliography}{99}

\bibitem{bracewell}
Bracewell, R. The Fourier Transform and Its Applications, 3rd ed. New York: McGraw-Hill, 1999.

\bibitem{sciencedirect-integral}
ScienceDirect (2024) \textit{Integral Transform}
Recuperado de \url{https://www.sciencedirect.com/topics/mathematics/integral-transform}

\bibitem{britannica-integral}
Britannica (2024) \textit{Integral Transform}
Recuperado de \url{https://www.britannica.com/science/integral-transform}

\bibitem{studysmarter-transformadas}
StudySmarter (2024) \textit{Transformadas Integrales}
Recuperado de \url{https://www.studysmarter.es/resumenes/ingenieria/ingenieria-electrica/transformadas-integrales/}

\bibitem{geeksforgeeks-fourier}
GeeksforGeeks (2024) \textit{Fourier Transform}
Recuperado de \url{https://www.geeksforgeeks.org/maths/fourier-transform/}

\bibitem{wolfram-dft}
Wolfram MathWorld (2024) \textit{Discrete Fourier Transform}
Recuperado de \url{https://mathworld.wolfram.com/DiscreteFourierTransform.html}

\bibitem{wolfram-convolution}
Wolfram MathWorld (2024) \textit{Convolution}
Recuperado de \url{https://mathworld.wolfram.com/Convolution.html}

\bibitem{youtube-convolution}
3Blue1Brown (2024) \textit{But what is a convolution?}
Recuperado de \url{https://www.youtube.com/watch?v=IaSGqQa5O-M}

\bibitem{statlect}
Taboga, Marco (2021). "Discrete Fourier transform", Lectures on matrix algebra. \url{https://www.statlect.com/matrix-algebra/discrete-Fourier-transform.}

 


\end{thebibliography}

\end{document}